\section{Experiments}
% Integration checklist for the final experiment pass.
% Confirmed inputs already reflected below:
% (1) DHC/TGN-edge/STHN/RepeatMixer/EdgeBank MRR tables from results/edge_only_formal/reports/.
% (2) Classical ML baselines from results/edge_only_formal/reports/tabular_ml_*.csv.
% (3) Existing ablation/sensitivity figures under ICONIP/figures/.
% (4) Existing cross-relation Patent-US figure under ICONIP-old/figures/.
% (5) Existing relation-memory/context case table under results/relation_temporal_context_case_study/.
% Pending final integration, do not describe as completed until files are present:
% (a) DyGFormer/TGAT successful model-zoo metrics, if reruns produce non-missing CSV rows.
% (b) Newly generated memory trajectory/convergence/context-heatmap/error-case figures.
% (c) Any final visual/case-study text from the baseline/visual/case agents.
% Page-budget note: keep only MRR main tables in the paper body unless AP/AUC tables replace a figure or move to appendix/supplement.
\subsection{Experiment Setup}
\noindent\textbf{Datasets.} We evaluate DHC on four processed temporal heterogeneous graphs from scholarly citations \cite{aminer,dblp}, a locally archived US-patent export, and the THGL software benchmark \cite{ahrabian2020software}. 
The scholarly graphs contain citation or coauthor dynamics. 
The patent graph contains organization--technology-code links with auxiliary patent relations. 
The software graph contains typed temporal interactions among software entities. Table~\ref{tab:dataset_stats} summarizes the processed statistics.
\rev{Patent-US was constructed from the archived \texttt{all-us.csv}, \texttt{IPCSubClass\_US.csv}, and \texttt{clsIII-reInd.csv} files. The export contains patent identifiers, publication dates, assignee names, IPC codes, and abstracts. We normalize assignee names, map IPC codes to level-III technology classes, and create patent--organization, patent--class, organization--class, and organization--organization events at the patent publication time. The prediction target is organization--class; the other three relations are context. The previously cited patent-statistics webpage was not the graph source and has therefore been removed. Because the upstream download URL/version is not recoverable from the archived project metadata, we do not claim a more specific provenance; a reproducibility release must include the archive checksum and redistribution URL or access instructions.}

\begin{table}[t]
\centering
\caption{Statistics of the processed temporal datasets. R-Types denotes relation types, N-Types denotes node types, and Avg.D denotes average degree.}
\label{tab:dataset_stats}
\scriptsize
\setlength{\tabcolsep}{5pt}
\begin{tabular}{lrrrr}
\toprule
Statistic & aminer-paper & dblp & patent-us & thgl-software \\
\midrule
Nodes       & 265,234 & 184,938 & 201,321 & 681,927 \\
Edges       & 961,282 & 200,000 & 1,047,764 & 1,489,806 \\
R-Types     & 4       & 2       & 4         & 14 \\
N-Types     & 3       & 2       & 3         & 4 \\
Avg.D       & 7.24    & 2.16    & 10.41     & 4.37 \\
\bottomrule
\end{tabular}
\end{table}

\noindent\textbf{Task views and data splits.}
All task views are built from the full processed event stream. 
A selected target relation provides supervised edges, while other relations remain as context edges. 
In transductive views, target edges are sorted by time and split into 70\%/15\%/15\% training, validation, and test parts. 
The exported windows keep all context edges before the corresponding cutoffs. Only target edges have \texttt{edge\_label=1}; context edges have \texttt{edge\_label=0} and update states only. 
\rev{For the AMiner-Paper and DBLP inductive views, source-side unseen nodes are sampled from the target relation with a fixed seed; their target edges are assigned to test and all training context touching them is removed. Patent-US instead uses a chronological organization cold-start split and removes held-out organizations from the training window. ``Unseen'' therefore means unseen in training event history; static features, types, and global IDs remain available. Importantly, the THGL-Software run reported here triggered the builder's insufficient-supervision fallback to a chronological temporal split. Its ``inductive'' column is not a strict unseen-node result and must not be interpreted as one. A genuine THGL cold-start rebuild and rerun are required before making a four-dataset inductive claim.}

\noindent\textbf{Compared baselines.}
We compare with completed temporal link-prediction baselines under the same processed views. 
TGN-edge augments the memory-based continuous-time graph baseline \cite{rossi2020tgn} with explicit edge-type features \cite{gastinger2024tgb2}. 
STHN is a compact continuous-time link-prediction baseline for dynamic interactions \cite{li2023sthn}. 
RepeatMixer models repeated interactions through repeat-aware sampling \cite{zou2024repeatmixer}. 
EdgeBank is a non-parametric historical-recurrence baseline. 
We also add tabular baselines using XGBoost, Random Forest, and Logistic Regression over hand-crafted pair features, including endpoint attributes, relation id, time gaps, historical pair counts, node degrees, relation frequency, and endpoint feature similarity. 
TemporalCount is a non-parametric forecasting baseline that scores a candidate edge from historical pair and relation-pair intensity. 
These methods cover memory-based, heterogeneous, repeat-aware, non-parametric history, and interpretable tabular strategies.
DyGFormer and TGAT are supported by the model-zoo runner, but their formal outputs are currently missing or failed for the reported tasks; therefore we reserve them for the final version only if successful reruns produce complete metrics.
% Final pass insertion point: add TAMI/TGN-base here only if they are selected for the compact main table or moved to a supplementary table.

\noindent\textbf{Evaluation metrics and implementation details.}
We use MRR as the main ranking metric in the paper body, and use AUC/AP in ablations and case studies when they clarify calibration or binary discrimination. 
MRR averages the reciprocal of the ranks of the first relevant item in a list of predicted answers. Training follows strict chronological order. At the beginning of each epoch, 
dynamic states are reset and the training stream is replayed.
 Validation and testing also use causal replay: 
 the model scores each target batch before updating states with the same batch. 
 Training uses one sampled negative destination per positive edge, 
 while evaluation uses 20 negatives. The default dimensions are 128 for node memory, 
 64 for relation embeddings, 32 for time encoding, 
 and 128 for message dimension. The best checkpoint is selected by validation MRR for edge tasks.
 We run all experiments on two NVIDIA RTX 3090 GPUs with 24GB RAM.
\rev{For node attributes, each processed dataset stores a 128-dimensional vector. AMiner-Paper and DBLP concatenate 64 standardized structural/type coordinates with 64 deterministic hybrid-hash text coordinates (word 1--2 grams and character 2--4 grams) derived from available titles, abstracts, venues, author names, years, and dataset metadata. Patent-US uses character 2--4-gram TF--IDF for 120 coordinates and eight structural coordinates, from patent abstracts, organization names, and IPC descriptions. THGL-Software has no natural-language field in our input; its 128-dimensional vector is a padded structural/type representation. No pretrained language model is fine-tuned.}

\rev{The reported $K=20$ protocol was selected as a common, computationally feasible sampled-ranking protocol across all four processed views and all rerun baselines; it is not the official dataset-specific TGB/TGB-2.0 negative sampler. Consequently, the tables support controlled comparison within our candidate sets, not a leaderboard-level state-of-the-art claim. We agree that MRR increases mechanically as the candidate set becomes easier/smaller. We have not fabricated unavailable $K=100$ or all-candidate results; these stricter evaluations, together with multi-seed variance, are required follow-up experiments.}

\rev{Efficiency follows directly from the bounded design. For batch size $b$, memory width $d_m$, relation width $d_r$, and at most $K_N$ stored neighbors per query node, recurrent updates cost $O(bd_m^2)$ and relation-conditioned projections cost $O(bd_md_r)$; local encoding costs $O(bK_Nd_m^2)$ up to attention constants. Persistent state is $O(|\mathcal V|d_m+|\mathcal R|d_m+|\mathcal V|K_N)$, in addition to model parameters. Relation memory itself adds only $O(|\mathcal R|d_m)$ state, whereas the two node-memory buffers dominate at scale. Saved run metadata records roughly 1.0--1.2 million trainable parameters across the four configurations. End-to-end runtime and peak-memory numbers were not measured under a single controlled hardware/software protocol, so we report no incomparable wall-clock figures and identify this as a limitation.}
\subsection{Performance Comparison}
\label{sec:performance_comparison}
As shown in Tables~\ref{tab:inductive_mrr} and~\ref{tab:transductive_mrr}, we compare our model with the baselines that have completed outputs on the four paper datasets.
The tables intentionally exclude candidate model-zoo baselines whose rows are still missing, and include classical ML baselines only where the current formal reports contain completed values.

\textbf{Inductive results.}
\rev{Table~\ref{tab:inductive_mrr} reports cold-start results for AMiner-Paper, DBLP, and Patent-US; the THGL-Software column retains the artifact's original ``inductive'' label but is a temporal-fallback result, as disclosed above.}
DHC is competitive on AMiner-paper, DBLP, and Patent-US, but the strengthened baseline set also shows where simple recurrence or compact temporal modeling is sufficient. 
On AMiner-paper and DBLP, DHC ranks second behind STHN and is clearly stronger than tabular and historical baselines, indicating that relation-aware temporal context helps when academic links depend on typed structure and text-derived node evidence. 
On Patent-US, tabular and count-based methods are unexpectedly strong because the target organization--technology relation contains repeated, highly regular interactions. 
On THGL-Software, STHN and RepeatMixer outperform DHC in the current run, suggesting that DHC's larger relation-memory machinery is not always beneficial when software interactions are sparse and relation-local recurrence dominates.

\begin{table}[t]
\centering
\caption{\rev{MRR for history-level cold-start temporal link prediction. THGL-Software is a temporal fallback, not a strict unseen-node split.} The best result is in bold, and the second best is underlined.}
\label{tab:inductive_mrr}
\scriptsize
\setlength{\tabcolsep}{4pt}
\resizebox{\textwidth}{!}{%
\begin{tabular}{lcccc}
\toprule
Method & AMiner-paper & DBLP & Patent-US & THGL-Software \\
\midrule
TGN-edge & 0.4239 & 0.5005 & \underline{0.9906} & 0.6296 \\
EdgeBank & 0.0476 & 0.0476 & 0.9726 & 0.0476 \\
XGBoost & 0.2646 & 0.2081 & \textbf{1.0000} & 0.1283 \\
Random Forest & 0.0905 & 0.0792 & \textbf{1.0000} & 0.0742 \\
Logistic Regression & 0.3570 & 0.2308 & \textbf{1.0000} & 0.5769 \\
TemporalCount & 0.0476 & 0.0476 & 0.9820 & 0.0476 \\
STHN & \textbf{0.8419} & \textbf{0.7842} & 0.9833 & \textbf{0.8936}\\
RepeatMixer & 0.2559 & 0.1513 & 0.7444 & \underline{0.8170} \\
DHC & \underline{0.7790} & \underline{0.7683} & \underline{0.9906} & 0.4703 \\
\bottomrule
\end{tabular}}
\end{table}
% Final pass insertion point: if DyGFormer/TGAT reruns succeed, add rows to Tables~\ref{tab:inductive_mrr} and~\ref{tab:transductive_mrr} and revise the ranking discussion.
% If using the refreshed six-dataset tabular-only report, add one compact sentence here instead of a new table unless the page budget allows it.
\textbf{Transductive results.} Table~\ref{tab:transductive_mrr} shows the transductive setting. 
This setting is usually easier than the inductive setting because training and test nodes may overlap. 
DHC remains strong on AMiner-paper and Patent-US, but it is not uniformly best after adding stronger baselines. 
STHN gives the best AMiner-paper and DBLP transductive MRR, while TGN-edge performs best on Patent-US and THGL-Software. 
The Patent-US table is especially informative: XGBoost and TemporalCount approach the best neural baseline, confirming that repeated organization--technology interactions make history intensity a powerful signal. 
In contrast, tabular baselines are much weaker on AMiner-paper, where text-aware node evidence and temporal structure must be combined. 
The DBLP transductive result also warns that DHC can overfit or mis-rank when the graph has few relation types and limited degree, so we treat this setting as a limitation rather than claiming uniform superiority.
\begin{table}[t]
\centering
\caption{MRR result for transductive temporal link prediction.The best result is in bold, and the second best is underlined.}
\label{tab:transductive_mrr}
\scriptsize
\setlength{\tabcolsep}{4pt}
\resizebox{\textwidth}{!}{%
\begin{tabular}{lcccc}
\toprule
Method & AMiner-paper & DBLP & Patent-US & THGL-Software \\
\midrule
TGN-edge & 0.3782 & 0.5724 & \textbf{0.9967} & \textbf{0.8913}\\
EdgeBank & 0.0476 & 0.0476 & 0.9446 & 0.3048\\
XGBoost & 0.2621 & 0.1691 & 0.9913 & 0.4400\\
Random Forest & 0.2457 & 0.0616 & 0.9681 & 0.3096\\
Logistic Regression & 0.3404 & \underline{0.6812} & 0.6659 & 0.3268\\
TemporalCount & 0.0476 & 0.0476 & \underline{0.9953} & 0.3120\\
STHN & \textbf{0.8419} & \textbf{0.7842} & 0.8459 & \underline{0.8760} \\
RepeatMixer & 0.5648 & 0.5106 & 0.5872 & 0.8199\\
DHC & \underline{0.7670} & 0.2548 & 0.9885 & 0.8000 \\
\bottomrule
\end{tabular}}
\end{table}

Overall, the expanded comparison gives a more precise picture. 
DHC is most useful when relation heterogeneity and node evidence must be combined with temporal context, as in academic inductive prediction. 
It is less advantageous when a dataset is dominated by simple repeated pairs or when compact temporal sequence models already capture the main signal.
\subsection{Ablation and Sensitivity}
% Final pass insertion point: replace or supplement Figures~\ref{fig:ablation}--\ref{fig:sensitivity} if the visual agent produces updated memory-trajectory, convergence, or attention/context plots.
Figure~\ref{fig:ablation} reports the contribution of the main DHC components on Patent-US and THGL-Software. Removing recent-neighbor context causes the largest degradation in most settings. This confirms that local temporal neighborhoods provide important evidence beyond isolated endpoint states. 
The effect of relation memory is more task dependent. Its removal has little influence on Patent-US transductive, where the target relation is dense and highly regular, 
\rev{but it reduces performance on Patent-US cold-start and on the THGL-Software view labeled inductive (which is a temporal fallback). The former supports the incomplete-history interpretation, whereas the latter only shows task-dependent temporal behavior.} Removing edge-type information also lowers performance, 
but the drop is not obvious especially on Patent-US transductive tasks where relation patterns are relatively stable, but becomes more visible in settings with richer heterogeneity. 
Overall, the full model gives the best or near-best result across all four splits. 
These results show that the three components play complementary roles: 
edge types provide explicit heterogeneous cues, 
relation memory tracks evolving relation semantics, 
and recent-neighbor context supplies local temporal evidence.

\begin{figure}[t]
\centering
\includegraphics[width=0.92\textwidth]{figures/3_ablation.png}
\caption{Results of ablation study on patent-us and thgl-software datasets.}
\label{fig:ablation}
\end{figure}
Figure~\ref{fig:sensitivity} shows the sensitivity of DHC to negative sampling and temporal-neighborhood construction on Patent-US. 
The four parameters affect the model in different ways. For \emph{typed negative ratio}, a moderate ratio gives stable gains, 
while a larger value leads to mixed effects. This suggests that type-consistent negatives help reduce reliance on coarse node-type cues, but too many such samples may reduce negative diversity. 
The effect of \emph{popular negative ratio} is more evident. 
A moderate value provides informative hard negatives, 
especially in the inductive setting, where the best performance is obtained at 0.1. 
The transductive setting continues to benefit as the ratio increases to 0.2, 
indicating that frequent destination nodes remain a major source of ranking confusion 
when node identities are already observed.

The preferred temporal window also differs across task settings. 
A larger \emph{recent neighbor size} is more helpful for inductive prediction, 
where broader local evidence can compensate for limited historical information about unseen nodes. 
In contrast, the transductive setting favors a smaller window, which provides sufficient context 
while avoiding unnecessary noise. 
Finally, \emph{neighbor aggregation} has the smallest effect among the four parameters. 
Attention-based aggregation is consistently better than mean pooling, 
although the margin is limited. This indicates that recent interactions in Patent-US are relatively coherent, 
while selective weighting still offers a safer default for relation-aware temporal modeling.

\begin{figure}[t]
\centering
\includegraphics[width=0.92\textwidth]{figures/4_sensitivity.png}
\caption{Sensitivity analysis on patent-us dataset.}
\label{fig:sensitivity}
\end{figure}
\subsection{Cross-relation Patent Analysis}
% Final pass insertion point: verify whether a newer cross-relation visual supersedes ICONIP-old/figures/cross_relation_mrr_base_vs_best.png.
Patent tasks expose a different behavior from academic and software graphs because the target relation is strongly tied to repeated organization--technology activity. 
Figure~\ref{fig:cross-relation-patent} summarizes the completed cross-relation Patent-US runs. 
The larger gains occur when auxiliary patent relations provide context that is temporally aligned with the target organization--code relation. 
When the auxiliary relation is weakly synchronized, adding more context gives smaller improvements and can even make the ranking depend on stale neighbors. 
This explains why simple historical baselines are strong in Tables~\ref{tab:inductive_mrr} and~\ref{tab:transductive_mrr}, while DHC still benefits when the relation-memory state captures which auxiliary relations are currently useful.

\begin{figure}[t]
\centering
\includegraphics[width=0.82\textwidth]{../ICONIP-old/figures/cross_relation_mrr_base_vs_best.png}
\caption{Cross-relation Patent-US analysis. The figure compares the base and best relation-context settings under the completed Patent-US cross-relation runs.}
\label{fig:cross-relation-patent}
\end{figure}

\subsection{Relation-memory and Temporal Context Case Study}
% Final pass insertion point: update this paragraph and the input table only after the case-study agent writes its final artifacts.
Table~\ref{tab:case-relation-temporal-context} gives two aggregate cases grounded in the saved run artifacts. 
We treat them as aggregate case studies rather than per-edge rank examples, because the available outputs contain \texttt{metrics.json}, \texttt{train\_history.csv}, and memory trace files, but not exported edge-level before/after ranking trajectories. 
\rev{For the software artifact named \texttt{thgl\_software\_edge\_inductive}---which the preprocessing audit identifies as a temporal fallback---the controlled \texttt{best\_config\_module\_ablation} runs show that removing relation memory changes MRR only from 0.7971 to 0.7967, but lowers AP from 0.5661 to 0.5466.}
This 0.0195 AP drop (3.45\% relative to the full model) indicates that relation memory improves aggregate positive-edge discrimination even when the top reciprocal-rank metric is nearly unchanged. 
Removing community/context information causes a larger AP drop of 0.0382, matching the multi-relation and sparse-history nature of the software graph.
For the patent temporal-context case, the \texttt{patent\_cn} recent-neighbor sweep shows that inductive prediction prefers a moderate window ($K=16$, AP 0.9870), while the transductive split peaks at a shorter window ($K=8$, AP 0.9873). 
Increasing the inductive window to $K=32$ lowers AP to 0.9822, and the weakest transductive setting is $K=12$ with AP 0.9825. 
This supports the interpretation that a short window can miss useful context for weakly observed nodes, whereas an overly long or poorly aligned window may reintroduce outdated organization--technology evidence.

\begin{table}[t]
\centering
\caption{Aggregate case study of relation memory and temporal context. Panel (a) reports controlled \texttt{best\_config\_module\_ablation} runs on \texttt{thgl\_software\_edge\_inductive}; $\Delta$AP is computed as Full minus the ablated variant, so a positive value means the removed module was helpful. Panel (b) reports the \texttt{best\_config\_sensitivity} temporal-context sweep on \texttt{patent\_cn}; $\star$ marks the best AP within each split. No per-edge rank-change sample is reported because the saved artifacts provide aggregate metrics, train histories, and memory traces rather than exported edge-level ranking trajectories.}
\label{tab:case-relation-temporal-context}
\centering
\small
\textbf{(a) Relation-memory ablation}\vspace{2pt}
\begin{tabular}{lrrrr}
\toprule
Variant & MRR & AUC & AP & $\Delta$AP \\
\midrule
full & 0.7971 & 0.9726 & 0.5661 & 0.0000 \\
no\_relation\_memory & 0.7967 & 0.9730 & 0.5466 & 0.0195 \\
no\_edge\_type & 0.7908 & 0.9724 & 0.5509 & 0.0152 \\
no\_community & 0.7896 & 0.9712 & 0.5280 & 0.0382 \\
\bottomrule
\end{tabular}

\vspace{6pt}
\centering
\small
\textbf{(b) Temporal-context window}\vspace{2pt}
\begin{tabular}{lrrrrr}
\toprule
Split & $K$ & MRR & AUC & AP & $\Delta$AP \\
\midrule
CN-inductive & 8 & 0.9954 & 0.9995 & 0.9840 & -0.0030 \\
CN-inductive & 12 & 0.9955 & 0.9995 & 0.9855 & -0.0015 \\
CN-inductive & 16 & 0.9962 & 0.9995 & 0.9870 $\star$ & 0.0000 \\
CN-inductive & 24 & 0.9945 & 0.9994 & 0.9833 & -0.0037 \\
CN-inductive & 32 & 0.9934 & 0.9993 & 0.9822 & -0.0048 \\
CN-transductive & 8 & 0.9957 & 0.9996 & 0.9873 $\star$ & 0.0000 \\
CN-transductive & 12 & 0.9956 & 0.9995 & 0.9825 & -0.0047 \\
CN-transductive & 16 & 0.9952 & 0.9995 & 0.9840 & -0.0032 \\
CN-transductive & 24 & 0.9957 & 0.9996 & 0.9867 & -0.0005 \\
CN-transductive & 32 & 0.9957 & 0.9996 & 0.9857 & -0.0016 \\
\bottomrule
\end{tabular}
\end{table}


